% Generated draft: requires mathematical and source review.
\documentclass[11pt]{article}
\usepackage[margin=25mm]{geometry}
\usepackage{fontspec}
\setmainfont{Latin Modern Roman}
\usepackage{amsmath,amssymb,mathtools,unicode-math}
\setmathfont{Latin Modern Math}
\usepackage{enumitem,needspace,xurl,hyperref}
\setlength{\parindent}{0pt}
\setlength{\parskip}{5pt}
\setlength{\emergencystretch}{3em}
\newcommand{\N}{\mathbb{N}}
\newcommand{\Z}{\mathbb{Z}}
\newcommand{\Q}{\mathbb{Q}}
\newcommand{\R}{\mathbb{R}}
\newcommand{\C}{\mathbb{C}}
\newcommand{\F}{\mathbb{F}}
\newcommand{\A}{\mathbb{A}}
\newcommand{\PP}{\mathbb{P}}
\newcommand{\E}{\mathbb{E}}
\newcommand{\eps}{\varepsilon}
\newcommand{\dd}{\,\mathrm{d}}
\begin{document}
\section*{UnsolvedMath ID: 6400006}
\textbf{Source problem code:} catalogue entry
\subsection{Definitions and mathematical statement}
For an integer $n\geq1$, consider a real-valued $C^2$ function $u$ on all of $\C^n$ with positive-definite complex Hessian
\[
u_{i\bar j}=\frac{\partial^2u}{\partial z_i\partial\bar z_j},\qquad
(u_{i\bar j})>0,\qquad \det(u_{i\bar j})=1.
\]
These are the classical strictly plurisubharmonic determinant-one complex Monge--Amp\`ere solutions. The source asks for an appropriate complex analogue of the J\"orgens--Calabi--Pogorelov classification of global real convex determinant-one solutions.

The requested complex classification is not specified further. Plurisubharmonicity is not the same as real convexity, so real convexity cannot be silently imported. Also, adding any real pluriharmonic function $v$ with $\partial\bar\partial v=0$ preserves the complex Hessian and equation. In particular the potential is not determined by its associated K\"ahler metric, and an assertion that every potential itself is quadratic would require normalization or an equivalence convention.

The record does not choose a regularity/weak-solution class, growth condition, completeness of the associated metric, finite-energy hypothesis, coordinate equivalence, or a proposed list of models. The displayed $C^2$ positive-Hessian class is the standard elliptic classical interpretation, not a claim that the source supplied all these hypotheses. One must specify which additional conditions constitute an ``appropriate'' analogue before a definite classification theorem can be stated. The summary's various restricted results do not establish an unrestricted theorem.
\subsection{Short English statement}
Classify entire strictly plurisubharmonic solutions of the determinant-one complex Monge--Amp\`ere equation in an appropriate analogue of the real J\"orgens--Calabi--Pogorelov theorem. Growth, completeness, and the potential's pluriharmonic ambiguity need explicit treatment.
\subsection{Sources}
\begin{itemize}
\item Local UnsolvedMath registry, permanent ID 6400006, exact source code catalogue entry, ``Complex J\"orgens--Calabi--Pogorelov Theorem''; proposed by Simon Donaldson, 2008. Registry status partially solved, classification PARTIAL-PROGRESS.
\item Simon Donaldson, \emph{Some Problems in Differential Geometry and Topology} (2008), explicit question/program 6, \url{https://wwwf.imperial.ac.uk/~skdona/NONLINEARITYPROBLEMS.PDF}. Supplied normalized pdf-text extraction accessed July 29, 2026. No PDF was accessed here.
\item The supplied research summary reports classifications under various growth, radial, or finite-energy hypotheses and no verified fully general complex classification. It gives no precise theorem, model list, or publication references, so those claims are retained as supplied limited-status assessment rather than premises fixing the target.
\item Provenance says the author survey presents the programme as outstanding and flags current status and statement-reuse terms NEEDS\_REVIEW; the document is publicly author-hosted.
\end{itemize}
\end{document}
